\subsection{SAVE-Area Layout Discovery}
Class \texttt{0x00000002}, leaf \texttt{0x0004} describes the memory image used by SAVE and RESTORE. Index zero is
the common (:term:base.terms.header:). Indexes one and two describe the complete image and its fixed block. Each available extension
component contributes descriptor A followed by descriptor B, so for $N$ components
(:cpuid-field-target:base.cpuid.IMPLEMENTATION.SAVE_AREA_LAYOUT.HEADER.MAX_INDEX:)\texttt{MAX\_INDEX}\texttt{=2+2N}. Component descriptors are ordered by increasing (:ref:base.cpuid.IMPLEMENTATION.SAVE_AREA_LAYOUT.COMPONENT_DESCRIPTOR_A.COMPONENT_ID:).

\Needspace{1.25in}
\BedrockTableCaption{SAVE-AREA-LAYOUT Indexes}
\begingroup\footnotesize
\setlength{\aboverulesep}{0pt}
\setlength{\belowrulesep}{0pt}
\setlength{\extrarowheight}{1.2pt}
\begin{center}
\begin{tabularx}{\linewidth}{@{}>{\raggedright\arraybackslash}p{0.80in}>{\raggedright\arraybackslash}X@{}}
\toprule
\rowcolor{BedrockHeaderFill}
\textbf{Index} & \textbf{Contents}\\
\midrule
0 & common CPUID leaf (:term:base.terms.header:)\\
1 & (:cpuid-field-target:base.cpuid.IMPLEMENTATION.SAVE_AREA_LAYOUT.SAVE_AREA_SIZE.SAVE_AREA_SIZE_BYTES:)\texttt{SAVE\_AREA\_SIZE\_BYTES} in bits 63..0\\
2 & (:cpuid-field-target:base.cpuid.IMPLEMENTATION.SAVE_AREA_LAYOUT.LAYOUT.FIXED_SIZE_BYTES:)\texttt{FIXED\_SIZE\_BYTES[15:0]}, (:cpuid-field-target:base.cpuid.IMPLEMENTATION.SAVE_AREA_LAYOUT.LAYOUT.COMPONENT_COUNT:)\texttt{COMPONENT\_COUNT[31:16]}, (:cpuid-field-target:base.cpuid.IMPLEMENTATION.SAVE_AREA_LAYOUT.LAYOUT.BITMAP_WORDS:)\texttt{BITMAP\_WORDS[47:32]}, (:cpuid-field-target:base.cpuid.IMPLEMENTATION.SAVE_AREA_LAYOUT.LAYOUT.FMT:)\texttt{FMT[51:48]}; bits 63..52 are zero\\
\texttt{3+2i} & component descriptor A for zero-based component ordinal \texttt{i}\\
\texttt{4+2i} & component descriptor B for zero-based component ordinal \texttt{i}\\
\bottomrule
\end{tabularx}
\end{center}
\endgroup

\begin{BedrockListedFormatDiagram}{SAVE-Area Layout CPUID Result Formats}
\BedrockFormatRowRange{index 2}{63}{0}{%
\BedrockFormatReserved{(:term:base.terms.reserved:)}{12}
\BedrockFormatField{FMT}{4}
\BedrockFormatField{BITMAP}{16}
\BedrockFormatField{COUNT}{16}
\BedrockFormatField{FIXED\_SIZE}{16}
}
\BedrockFormatRowRange{descriptor A}{63}{0}{%
\BedrockFormatField{OFFSET}{32}
\BedrockFormatReserved{(:term:base.terms.reserved:)}{10}
\BedrockFormatField{BIT}{6}
\BedrockFormatField{ID}{16}
}
\BedrockFormatRowRange{descriptor B}{63}{0}{%
\BedrockFormatReserved{(:term:base.terms.reserved:)}{8}
\BedrockFormatField{INIT}{8}
\BedrockFormatField{ALIGN}{16}
\BedrockFormatField{MAX\_SIZE}{32}
}
\end{BedrockListedFormatDiagram}
{\small Diagram labels abbreviate the (:term:base.terms.field:) names in the descriptor tables below; \texttt{BITMAP}, \texttt{COUNT}, and \texttt{FIXED\_SIZE} are (:ref:base.cpuid.IMPLEMENTATION.SAVE_AREA_LAYOUT.LAYOUT.BITMAP_WORDS:), (:ref:base.cpuid.IMPLEMENTATION.SAVE_AREA_LAYOUT.LAYOUT.COMPONENT_COUNT:), and (:ref:base.cpuid.IMPLEMENTATION.SAVE_AREA_LAYOUT.LAYOUT.FIXED_SIZE_BYTES:).\par}

\Needspace{2.5in}
\Needspace{1.25in}
\BedrockTableCaption{SAVE Component Descriptor A}
\begingroup\footnotesize
\setlength{\aboverulesep}{0pt}
\setlength{\belowrulesep}{0pt}
\setlength{\extrarowheight}{1.2pt}
\begin{center}
\begin{tabularx}{\linewidth}{@{}>{\raggedright\arraybackslash}p{0.65in}>{\raggedright\arraybackslash}p{1.70in}>{\raggedright\arraybackslash}X@{}}
\toprule
\rowcolor{BedrockHeaderFill}
\textbf{Bits} & \textbf{Field} & \textbf{Meaning}\\
\midrule
\texttt{15..0} & (:cpuid-field-target:base.cpuid.IMPLEMENTATION.SAVE_AREA_LAYOUT.COMPONENT_DESCRIPTOR_A.COMPONENT_ID:)\texttt{COMPONENT\_ID} & stable component identifier\\
\texttt{21..16} & (:cpuid-field-target:base.cpuid.IMPLEMENTATION.SAVE_AREA_LAYOUT.COMPONENT_DESCRIPTOR_A.BITMAP_BIT:)\texttt{BITMAP\_BIT} & bit selecting this component\\
\texttt{31..22} & (:term:base.terms.reserved:) & zero\\
\texttt{63..32} & (:cpuid-field-target:base.cpuid.IMPLEMENTATION.SAVE_AREA_LAYOUT.COMPONENT_DESCRIPTOR_A.OFFSET_BYTES:)\texttt{OFFSET\_BYTES} & offset from the save-area base\\
\bottomrule
\end{tabularx}
\end{center}
\endgroup


\Needspace{1.25in}
\BedrockTableCaption{SAVE Component Descriptor B}
\begingroup\footnotesize
\setlength{\aboverulesep}{0pt}
\setlength{\belowrulesep}{0pt}
\setlength{\extrarowheight}{1.2pt}
\begin{center}
\begin{tabularx}{\linewidth}{@{}>{\raggedright\arraybackslash}p{0.65in}>{\raggedright\arraybackslash}p{1.70in}>{\raggedright\arraybackslash}X@{}}
\toprule
\rowcolor{BedrockHeaderFill}
\textbf{Bits} & \textbf{Field} & \textbf{Meaning}\\
\midrule
\texttt{31..0} & (:cpuid-field-target:base.cpuid.IMPLEMENTATION.SAVE_AREA_LAYOUT.COMPONENT_DESCRIPTOR_B.MAX_SIZE_BYTES:)\texttt{MAX\_SIZE\_BYTES} & allocated component size\\
\texttt{47..32} & (:cpuid-field-target:base.cpuid.IMPLEMENTATION.SAVE_AREA_LAYOUT.COMPONENT_DESCRIPTOR_B.ALIGNMENT_BYTES:)\texttt{ALIGNMENT\_BYTES} & required component alignment\\
\texttt{55..48} & (:cpuid-field-target:base.cpuid.IMPLEMENTATION.SAVE_AREA_LAYOUT.COMPONENT_DESCRIPTOR_B.INIT_POLICY:)\texttt{INIT\_POLICY} & 1 restores the component-defined initial state when its bitmap bit is clear\\
\texttt{63..56} & (:term:base.terms.reserved:) & zero\\
\bottomrule
\end{tabularx}
\end{center}
\endgroup


(:ref:base.cpuid.IMPLEMENTATION.SAVE_AREA_LAYOUT.SAVE_AREA_SIZE.SAVE_AREA_SIZE_BYTES:) is the complete byte range validated and accessed by SAVE or RESTORE. It is at
least (:ref:base.cpuid.IMPLEMENTATION.SAVE_AREA_LAYOUT.LAYOUT.FIXED_SIZE_BYTES:), covers the end of every component, and is a multiple of 64 bytes.
(:ref:base.cpuid.IMPLEMENTATION.SAVE_AREA_LAYOUT.LAYOUT.FIXED_SIZE_BYTES:) is \texttt{0x00c0}, (:ref:base.cpuid.IMPLEMENTATION.SAVE_AREA_LAYOUT.LAYOUT.BITMAP_WORDS:) is 1, and (:ref:base.cpuid.IMPLEMENTATION.SAVE_AREA_LAYOUT.LAYOUT.FMT:) is zero for the
save-image format defined by this specification.

Each descriptor names one available extension component. Component IDs and bitmap bits are unique. Component
regions begin at or after the fixed block, meet their reported alignment, and do not overlap. A set state-block
bitmap bit selects the descriptor carrying the same (:ref:base.cpuid.IMPLEMENTATION.SAVE_AREA_LAYOUT.COMPONENT_DESCRIPTOR_A.BITMAP_BIT:). (:ref:base.cpuid.IMPLEMENTATION.SAVE_AREA_LAYOUT.COMPONENT_DESCRIPTOR_B.INIT_POLICY:)\texttt{=1} means that
RESTORE reconstructs the component-defined initial state when that component\textquotesingle{}s bitmap bit is clear.
